% Drawing 71/02-5 -- folder 71, batch 02, archive page 29 (PDF page 30).
% Transcription: page de croquis, sans ordre, autour du texte écrit en travers. À l'encre noire : en haut, un polygone en pointe de flèche dont un creux est fermé en tirets, sous un trait isolé ; au milieu, deux polygones en dents de scie dont les creux sont fermés en pointillé ; plus bas, deux quadrilatères contenant chacun une étoile à quatre branches en tirets, appuyée sur leurs sommets ; en dessous, un petit triangle $s$, $s'$, $s''$ avec une ligne pointillée de $s$ vers un point $t$, et un polygone non convexe portant $s$, $s'$, $s''$ et $t$, une verticale en tirets descendant de $s$ vers $t$. Au crayon, pâles : de longs triangles effilés, trois petits triangles surchargés en haut à droite, un polygone en dents de scie et un triangle hachuré en bas.
% Kind: page of sketches: non-convex polygons with their hulls dotted, stars inside quadrilaterals. Method 2 (schematic redraw from the facsimile). Status: partial.
% Notes: Only the seven ink figures are redrawn, each closed through its vertices as read: (1) top, the arrowhead polygon with the dashed stroke closing its notch, and the isolated stroke above it; (2) middle left, a saw-tooth polygon with a dotted stroke over its upper notch and a dashed one closing its lower notch, the horizontal side running on a little to the right; (3) middle, the larger saw-tooth polygon, its upper left notch closed by a dotted stroke ending at a heavy point, a second heavy point at the next tooth, the bottom closed by a dashed stroke; (4), (5) the two quadrilaterals, each with a dashed four-pointed star on its four vertices; (6) the small triangle s, s', s'' (two heavy sides from s) with the dotted line from s to the mark read t; (7) the non-convex polygon with s at the tip of its notch, s' and s'' at the foot of the notch, joined by a dotted stroke, and the dashed vertical from s down to the point t. Left out: all the pencil sketches (the long thin triangles, the three small overworked triangles at the upper right, the saw-tooth polygon and the hatched triangle at the bottom, a pencil label near figure 3), a small arrow-shaped polygon drawn inside the upper teeth of figure 3, and the text written across the sheet. Still differs: the teeth at the top of figure 3 and the inner vertices of the two stars are approximate; in figure 5 a small zigzag near the top vertex is not drawn. Decided by a machine pass, 2026-11-12: the mark at the end of the dotted line of figure 6 is read t, as the letter at the foot of figure 7 (same shape; the text across the page speaks of s, s', s'' and t). Not checked by a reader. The transcription's note mentions two drawings only (the pencil saw-tooth polygon in a large triangle and the ink arrowhead polygon); the page carries about a dozen. Redrawn 2026-11-12 (issue 128): to check again. 2026-11-13, Claude Opus 5.5, overlay check (issue 147): the whole sheet laid over the scan rendered at 350 dpi in colour; not compiled (no TeX engine in the session), the TikZ was rasterised by a script; nothing changed. The seven ink figures lie on their ink within about 1 to 2 units, apart from the upper teeth of figure 3 and the inner vertices of the two stars, as said above. Kept partial because the mark names the pale pencil sketches too (the long thin triangles, the three overworked small triangles at the upper right, the pencil saw-tooth polygon and the hatched triangle at the bottom), and none of them is drawn; they are faint, overworked and partly under the writing, and would have to be redrawn stroke by stroke.
% Coordinates: 1 unit = 2 mm = 1/100 of the crop width (bbox in page fractions: [0.18, 0.07, 0.96, 0.69]).
% Redrawn by Claude (claude-fable-5-2), issue #128, 2026-11-12; not checked by a human; not compiled (no TeX in the session).
% Rights: reproduces part of the work -- behind RENDER_TRANSCRIPTS, never published by itself.
\documentclass[tikz,border=3pt]{standalone}
\usepackage{amsmath,amssymb}
\usetikzlibrary{patterns,arrows.meta}
\tikzset{pencil/.style={black!45}, struck/.style={black!35}, mpt/.style={circle,draw,inner sep=1.3pt}, pt/.style={circle,fill,inner sep=0.9pt}, opt/.style={circle,draw,fill=white,inner sep=1.1pt}}
\begin{document}
\begin{tikzpicture}[x=2mm,y=2mm,line cap=round,line join=round,font=\scriptsize]
\draw (32.3,-12.9) -- (40.8,-7.1);
\draw (33.9,-15) -- (40.3,-15.3) -- (43,-10) -- (44.1,-23) -- cycle;
\draw[dashed] (35.2,-14.1) -- (43,-10);
\draw (12.8,-30.5) -- (15.6,-32.7) -- (16,-28.8) -- (22.7,-34.9) -- (15,-34.7) -- (12.9,-40.5) -- cycle;
\draw (22.7,-34.9) -- (25,-35);
\draw[densely dotted] (12.8,-30.5) -- (16,-28.8);
\draw[dashed] (12.9,-40.5) -- (22.7,-34.9);
\draw (40.7,-31) -- (43.3,-29.7) -- (45.1,-32.8) -- (45.5,-28.8) -- (47,-27.9) -- (47.8,-31.4) -- (48.5,-26.9) -- (50.1,-26.2) -- (60.4,-37.9) -- (48.8,-35.4) -- (47.8,-39.3) -- (44.2,-36.1) -- (41.3,-40.7) -- cycle;
\draw[densely dotted] (43.3,-29.7) -- (45.5,-28.8);
\node[pt] at (45.5,-28.8) {};
\node[pt] at (47,-27.9) {};
\draw[dashed] (41.3,-40.7) -- (59.3,-38.1);
\draw (34,-56.1) -- (42.1,-52) -- (52.9,-60.9) -- (34.3,-63.3) -- cycle;
\draw[dashed] (34,-56.1) -- (39.5,-56.4) -- (42.1,-52);
\draw[dashed] (42.1,-52) -- (44.2,-57.3) -- (52.9,-60.9);
\draw[dashed] (52.9,-60.9) -- (40.1,-58.9) -- (34.3,-63.3);
\draw[dashed] (34.3,-63.3) -- (37,-58.3) -- (34,-56.1);
\draw (63.8,-53.9) -- (69,-50.7) -- (82.5,-57.4) -- (67.4,-61.1) -- cycle;
\draw[dashed] (63.8,-53.9) -- (68.9,-53.8) -- (69,-50.7);
\draw[dashed] (69,-50.7) -- (72.9,-55.2) -- (82.5,-57.4);
\draw[dashed] (82.5,-57.4) -- (71.8,-57) -- (67.4,-61.1);
\draw[dashed] (67.4,-61.1) -- (68,-56.4) -- (63.8,-53.9);
\draw[very thick] (32.3,-77) -- (33.9,-81.7);
\draw[very thick] (32.3,-77) -- (36.9,-79.7);
\draw[densely dotted] (33.9,-81.7) -- (36.9,-79.7);
\draw[densely dotted] (32.3,-77) -- (42.2,-70.7);
\node[inner sep=1pt] at (30.5,-77.1) {\small $s$};
\node[inner sep=1pt] at (39,-78.6) {\small $s''$};
\node[inner sep=1pt] at (32.9,-83.4) {\small $s'$};
\node[inner sep=1pt] at (43.3,-69.5) {\small $t$};
\draw (53.1,-71.8) -- (63.4,-69.1) -- (71.1,-80.8) -- (56.4,-89) -- (63.8,-80) -- (62.7,-75.7) -- (56.8,-82.1) -- (54.4,-75.7) -- (51.9,-83.6) -- cycle;
\draw[dashed] (54.4,-75.7) -- (54.8,-89.8);
\draw[densely dotted] (51.9,-83.6) -- (56.8,-82.1);
\node[inner sep=1pt] at (55.4,-74.3) {\small $s$};
\node[inner sep=1pt] at (58.7,-81.8) {\small $s''$};
\node[inner sep=1pt] at (50.6,-85.8) {\small $s'$};
\node[inner sep=1pt] at (56.9,-90.6) {\small $t$};
\end{tikzpicture}
\end{document}
